\section{Reevaluation}

Each evaluation DAG represents an orbit's history which is valid with regards to
the initial evaluation. %
When reevaluating, editing the parametric specification makes the topological
parameters subject to changes. %
Thus, it is necessary to build reevaluation DAGs from the evaluation DAGs in
order to update topological parameters. %
Once built, a reevaluation DAG can designate one, several, or no orbit depending
on the editing of the parametric specification. %

In this section, we keep using the previous example (Fig.~\ref{fig:evaluation}
and \ref{fig:reevaluation}) and its edited parametric specification
(Fig.~\ref{fig:init-spec} and \ref{fig:reval-spec}), containing an added vertex
insertion on the square's edge right after its creation and the deletion of the
vertex insertion on the cube's edge. %

\subsection{Reevaluation DAG}

\begin{figure}[ht]
    \centering
    \includegraphics[height=.32\textheight]{pn3-reval}
    \caption{$PN_3$'s reevaluation DAG}
    \label{fig:pn3-reval-dag}
\end{figure}

Contrary to the related evaluation DAG, a reevaluation DAG is built top-down
throughout the reevaluation process. %
While an evaluation DAG represents the orbit's history of a topological
parameter, the reevaluation DAG represents the history of this very same orbit
after editing the parametric specification.

For example, consider $PN_{3}$ the topological parameter of \texttt{4-extrusion}
which extrudes the cube's top face into a second cube. %
The reevaluation process builds $PN_{3}$'s reevaluation DAG step-by-step from
its evaluation DAG %
% The $PN_3$ reevaluation DAG (Fig.~\ref{fig:pn3-reval-dag}) is built step-by-step
% during the reevaluation process from its evaluation DAG
(Fig.~\ref{fig:eval-dag-pn3}). %

% \begin{description}
% \item[Level 1]
\textbf{Level 1} The application \texttt{1-square} has no topological parameter
and, thus, is identically reevaluated. Once again, its node $n5$ creates a
single dart $5$ and the volume face $\orbit{0,1}(5)$ is identically reevaluated.
This is why the orbit level 1 of the reevaluation DAG references dart $5$.

        % \item[Level 2]
\textbf{Level 2} The second application is an added one. \texttt{ADD1-insert}
does match an edge of the tracked volume face $\orbit{0,1}(5)$ and modifies it,
as deduced from the rule. Therefore, the event level 2 contains
\texttt{MODIFICATION} and the orbit level 2 contains the same dart $5$ and orbit
$\orbit{0,1}$.

        % \item[Level 3]
\textbf{Level 3} The third application is the extrusion of the square face into a
cube. During its application, dart $5$ is matched by the hook $n_1$
(Fig.~\ref{fig:rule-face-extrusion}). From the evaluation DAG, the tracked dart
is the copy of dart $5$ created by node $n_6$. Applying the rule allows finding
out this dart, numbered $59$.

  % \item[Level 4]
\textbf{Level 4} Finally, the last application of $PN_3$ is deleted.
Consequently, the modification that occurred during the initial evaluation does
not occur at reevaluation. Therefore, the event level 4 contains
\texttt{NOMODIF} and the orbit level 4 contains the same dart~$59$.
% \end{description}

The reevaluation DAG identifies the concrete name using its persistent name. %
$PN_3$'s concrete name is $59$ (cf. Fig.~\ref{fig:2-reeval} and
\ref{fig:4-reeval}). %

\begin{figure}[ht]
    \centering
    \includegraphics[height=.45\textheight]{pn6-reval}
    \caption{$PN_6$ reevaluation DAG}
    \label{fig:pn6-reval-dag}
\end{figure}

We now study a more complex example with the reevaluation DAG of $PN_{6}$
(Fig.~\ref{fig:pn6-reval-dag}), which represents the pyramid's top vertex
(Fig.~\ref{fig:6-eval}). %

% \begin{description}
% \item[Level 1 to 4]
\textbf{Level 1 to 4} These three levels are similar to $PN_3$'s reevaluation
DAG (Fig.~\ref{fig:pn3-reval-dag}), with the tracking of the volume face vertex
$\orbit{1}$-orbit in addition to the $\orbit{0,1}$-orbit one. %
The $\orbit{1}(5)$-orbit is created by the \texttt{1-square} application. %
\texttt{ADD1-insert} matches and preserves one dart of this orbit with the node
$n_{0}$ of the insertion rule (Fig.~\ref{fig:rule-insert-vertex}). Thus, the
event level 2 contains \texttt{NOMODIF}. Then, this orbit is copied to create a
new volume face vertex by \texttt{2-extrusion}. Again,
\texttt{DELETE 3-insert} does not modify the orbits.

% \item[Level 5]
\textbf{Level 5} The fifth application \texttt{4-extrusion} of the extrusion
rule (Fig.~\ref{fig:rule-face-extrusion}) matches the volume face $\orbit{0,1}$
with its hook $n1$ and creates a copy of dart $59$ dart $63$. As in the initial
evaluation, it creates a volume vertex and a face. Then, the event level
contains two \texttt{CREATION}. The orbit level references dart $63$ and
contains both orbits $\orbit{1,2}$ and $@1.@2.\orbit{0,1}$.

% \item[Level 6]
\textbf{Level 6} Finally, while the tracking of \texttt{5-collapse} shows that
the application keeps merging the volume vertices incident to the matched face,
it preserves dart $63$ which is matched by node $n2$. Therefore, the event level
contains one \texttt{MERGE} and the orbit level references dart $63$ and
contains the volume vertex $\orbit{1,2}$.
% \end{description}

At last, $PN_6$'s concrete name is dart $63$ (cf Fig.~\ref{fig:6-reeval}). %

These two examples here are quite straightforward as there was only one possible
candidate dart each time. %
However, in some complex specifications, there can be more than a single dart to
choose between. %

\subsection{Parameter matching strategies}

The editing of the parametric specification leads to having a different DAG at
reevaluation (with event levels and/or branches being added, deleted or both). %
For example, an orbit split present in the evaluation DAG may disappear during
the reevaluation, a merging can be added and so on. %
Several matching strategies can then be considered depending on the
application's context. %
This can be illustrated with $PN_7$'s example which designates the face that
must be coloured (Fig.~\ref{fig:8-eval}). %
The addition of \texttt{ADD1-insert} application at reevaluation splits the
origin of designated face, resulting in the addition of a branch in the
reevaluation DAG. %


\begin{figure}[ht]
    \centering
    \includegraphics[width=\columnwidth]{pn7-reval}
    \caption{$PN_7$ reevaluation DAG}
    \label{fig:pn7-reval-dag}
\end{figure}

Let's work through $PN_7$'s reevaluation DAG shown in Fig.~\ref{fig:pn7-reval-dag}. %

% \begin{description}
% \item[Level 1]
\textbf{Level 1} As seen previously, \texttt{1-square} creates the tracked
orbits traced in the evaluation DAG.

% \item[Level 2] \textbf{Level 2} The added application
\textbf{Level 2} \texttt{ADD1-insert} matches dart $3$. Since the vertex insertion rule
(Fig.~\ref{fig:rule-insert-vertex}) splits the volume face edge $\orbit{0}(3)$,
its history is also split and there are two concrete names to consider. It
follows that the event level contains two \texttt{SPLIT}, one for each volume
face edge suborbit, and two \texttt{NOMODIF}, one for each dart $3$ and $4$.

% \item[Level 3] \textbf{Level 3} The third application
\textbf{Level 3} \texttt{2-extrusion} extrudes the square face into a cube. The extrusion rule
(Fig.~\ref{fig:rule-face-extrusion}) matches dart $3$ with its hook $n_1$. It
extrudes the face edge $\orbit{0}(3)$ into a volume face, the dart $\orbit{}(3)$
into a face edge and its node $n_2$ creates the dart $37$ as a copy of dart $3$.
The same goes for the history on the right. The rule matches dart $4$, it
creates the same two orbits and a dart $52$ as a copy of dart $4$. % Therefore,
both event levels contain two \texttt{CREATION}. Both orbit levels contain
orbits $\orbit{0,1}$ and $\orbit{0}$. %

% \item[Level 4] \textbf{Level 4} The last application,
\textbf{Level 4} \texttt{5-triangulation} triangulates the faces designated by the level above.
Both darts $37$ and $52$ are matched and preserved by the rule.
% \end{description}

Upon reevaluation, $PN_7$'s DAG matches to different darts. An option would be to colour only one face (either $\orbit{0,1}(37)$ or $\orbit{0,1}(52)$). %
Another option would be to apply \texttt{8-colour} two times, one for dart $37$ and one for dart $52$ (as shown in Fig.~\ref{fig:8-reeval} which represent our default strategy). %
In case the reevaluation DAG has two or more leaves, it shows all the possible entities that can be matched for a specific persistent name and to set a strategy up. %
Either way, such a strategy allows users to have a choice and best fit their modeling intents, depending on the application's context. %

%%% Local Variables:
%%% mode: latex
%%% TeX-master: "../main"
%%% End:
